\subsection*{从边值问题理解电容}

电容的定义看上去只是一个比值，但它实际上概括了完整的静电边值问题。以孤立导体为例，给定导体表面的
电势以后，导体外部的电势必须满足 Laplace 方程，并在导体表面取给定常数、在无穷远趋于零。唯一性定理
保证满足这些条件的解只有一个。若把表面电势由 $\varphi$ 改成 $k\varphi$，原解整体乘以 $k$ 就满足新的
边界条件；于是电场、表面电荷密度和总电荷也全部乘以 $k$。这正是 $Q=C\varphi$ 成立的根本原因。

在这个构造中，其他接地导体虽然电势为零，却不等于“没有作用”。它们表面会出现感应电荷并改变空间电场，
因而也会改变未接地导体的电容。若未接地导体带正电，邻近接地导体朝向它的一侧通常出现负感应电荷。
接地意味着大地可以供给或吸收电子，所以这些导体的净电荷不必保持为零。把接地导体移近，通常会使给定
电荷造成的电势降低，从而使 $C=Q/\varphi$ 增大。这说明“一个导体的电容”仍然包含周围边界的信息。

双极电容器也由同样的线性原理决定。两极带等量异号电荷时，正极发出的电场线最终终止在负极，理想情况下
电场主要局限在两极之间。电容器储存的不是“正电荷”或“负电荷”本身，而是把正负电荷分离后形成的电场能。
如果保持两极形状不变，只把它们靠近，给定 $Q$ 造成的电势差会减小，因此电容增大；如果增大正对面积，
给定总电荷分摊到更大表面上，场强与电势差降低，电容同样增大。平行板公式中的 $S$ 和 $d$ 依赖关系由此
获得直观解释，而不只是代数结果。

平行板近似还包含一个容易被忽略的尺度条件。若极板边长与间距同量级，边缘弯曲的电场占据很大比例，不能
把全部场区都视为匀强场；此时 $C=\varepsilon_0S/d$ 只能给出粗略数量级。将极板边长同时放大而保持间距
不变，边缘区域的宽度仍约为 $d$，但中部匀强区域所占比例越来越大，近似才逐渐变好。竞赛题若强调
$l\gg d$，就是在要求忽略这部分端部修正。

圆柱电容的积分也展示了几何如何进入电容。圆柱 Gauss 面的面积正比于 $r$，所以场强按 $1/r$ 衰减；
电势差是 $1/r$ 从 $a$ 到 $b$ 的积分，因而出现 $\ln(b/a)$。如果只记结论而忽略这一过程，容易把它与球形
电容的 $1/R_1-1/R_2$ 混淆。三种典型电容器可用统一步骤处理：确定对称性，选取 Gauss 面求 $E$，沿场线
积分求 $U$，最后用 $C=Q/U$ 消去 $Q$。

\subsection*{连接、充电和储能的完整推理}

并联电容的关键是端点相同。达到静电平衡后，理想导线上的电势处处相同，因此每只电容的两端电压都等于
外部端口电压。电源向整个并联系统搬运的电荷是各电容正极板电荷之和，故等效电容相加。串联电容的关键
则是中间节点孤立：这个节点由相邻两块极板和连接导线构成，外部没有导电通路。它可以在内部重新分配正负
电荷，却不能改变净电荷；初始中性时，相邻两电容上的电荷量绝对值必然相同。把每段电压 $Q/C_i$ 相加，
便得到倒数相加的规律。

这种推理比直接背公式更可靠。例如串联电容的中间节点若预先带有净电荷，各电容所带电荷量便不再简单相等；
若电容器之间不是理想导线而连接了别的元件，节点条件也要重新分析。所谓“串联”或“并联”并不是由图形
看起来像一条线还是两条支路决定，而是由电势约束和电荷守恒决定。

电容储能中的因子 $1/2$ 来源于充电过程中电压从零逐渐增加。若错误地把最终电压 $U$ 乘以全部电荷 $Q$，
会得到 $QU$，比真实储能大一倍。正确做法是把每一小份电荷在当时电压 $u=q/C$ 下搬运所需的功积分起来。
也可以在 $q-u$ 图上理解：充电过程对应从原点到 $(Q,U)$ 的直线，储能等于直线下的三角形面积
$QU/2$。对于线性介质，$Q$ 与 $U$ 成正比；若介质具有强非线性或滞后，储能必须沿真实充电路径积分，
不能直接套用 $CU^2/2$。

把电场能密度在空间积分，应当得到与充电做功相同的总能量。以平行板电容为例，板间场近似均匀，
\[
 \int u_E\dd V=\frac12\varepsilon_0E^2Sd
 =\frac12\frac{\varepsilon_0S}{d}(Ed)^2
 =\frac12CU^2.
\]
这一等价性说明，集中参数 $C$ 的储能公式与分布在空间中的场能描述的是同一个物理事实。遇到介质只填充
部分空间、场分布不均匀的情形，直接积分场能有时比先求等效电容更方便。

电源给电容充电时还要区分“电源做功”和“电容最终储能”。理想恒压源每搬运 $\dd q$ 都做功
$U_0\dd q$，所以总功为 $U_0Q$；而电容电压从零升到 $U_0$，平均电压只有 $U_0/2$，最终储能是
$U_0Q/2$。差额必须进入电阻发热、辐射或电磁场的暂态过程。把电阻取得很小只会让过程更快，不会凭空改变
初态与终态之间的能量差。

\subsection*{倒易定理的三种典型用法}

Green 倒易定理的力量在于，它允许把一个难求的响应换成另一个容易计算的响应。应用时必须保证两个状态中
导体的形状、位置和介质完全相同，仅改变各导体的电荷或电势。若几何也改变，两个状态不再属于同一个线性
边值问题，不能直接使用倒易关系。

第一种典型用法是交换“点电荷”和“测量电势的位置”。在导体球外放点电荷时，球面感应电荷分布很复杂；
但把球改成带单位电荷、去掉外部点电荷后，球面电荷均匀，任意外点电势都容易写出。倒易定理把第二种状态
的已知电势转换成第一种状态中整个球的电势，从而避开求感应电荷密度的繁重计算。

第二种用法是交换两个导体。状态 I 中让导体 1 带电 $Q$、导体 2 中性；状态 II 中让导体 2 带相同电荷
$Q$、导体 1 中性。代入倒易式可知，第一状态中导体 2 的电势等于第二状态中导体 1 的电势。这个结论与
两导体的形状是否相同、距离多远都无关，只要求几何在交换状态时没有改变。

第三种用法涉及接地感应电荷。状态 I 中把导体 1 提升到单位电势而让导体 2 接地，状态 II 中交换两导体。
倒易定理说明，两次实验中接地导体上出现的总感应电荷相同。用矩阵语言表述，就是电容系数的互易性
$C_{12}=C_{21}$。这一性质使复杂导体系统所需的独立系数数量大幅减少，也是后来网络互易定理的静电学原型。

倒易定理还可以和能量观点联系。导体组总静电能可写成
\[
 W=\frac12\sum_iQ_i\varphi_i
   =\frac12\sum_{ij}C_{ij}\varphi_i\varphi_j.
\]
交换两个电势变量不应改变同一个能量函数的混合二阶变化，因此系数矩阵必须对称。这给出了倒易关系的另一
种直观来源：互易性并非偶然的代数巧合，而是线性系统存在单值能量函数的结果。

\subsection*{从实验现象到介质极化}

插入电介质的实验可以有两种做法。若电容器先充电后与电源断开，自由电荷 $Q$ 不变，插入介质后电压表读数
降低；由 $C=Q/U$ 可知电容增大。若电容器始终与恒压电源相连，电压不变，插入介质时电源继续向极板输送
电荷；由 $Q=CU$ 可知储荷量增加。两种实验观测的是同一事实，只是外部保持不变的量不同。

同一种介质填满不同形状的电容器时，电容近似按同一比例增加，这一点非常关键。它说明介质作用不能只用
某种特定极板几何来解释，而应由材料自身的响应参数描述。相对介电常数正是这样的宏观参数。表 32.1 中
不同物质的数值相差很大，水的相对介电常数远大于干燥空气，表明分子结构和极化机制会显著影响宏观响应。
击穿场强则描述另一件事：介质在多强的场下还能维持绝缘。介电常数大并不意味着击穿场强必然大，设计真实
电容器时必须同时考虑储能能力和绝缘可靠性。

在保持自由电荷不变的实验中，插入介质后宏观场减小。微观上，外场使原子或分子的正负电荷中心分离，
形成与外场同向的电偶极矩；大量偶极子在介质两端留下束缚电荷。束缚电荷产生的场在介质内部通常与原外场
反向，于是合场减弱。这里“偶极矩与外场同向”和“偶极子附加场与外场反向”并不矛盾：前者描述单个偶极
矩的方向，后者描述极化介质端面束缚电荷在材料内部产生的退极化场。

均匀极化时介质内部的束缚体电荷为零，但表面束缚电荷一般不为零；非均匀极化时 $\rho_b=-\nabla\cdot\vect P$
也会贡献电场。因此，不能把极化简单理解为“只有表面有电荷”。具体是表面项还是体积项占主导，取决于
$\vect P$ 的空间分布和边界形状。

引入 $\vect D$ 的主要好处，是在已知自由电荷而束缚电荷难以逐个计算时直接使用 Gauss 定律。例如平行板
间完全填充均匀线性介质，极板自由电荷面密度为 $\sigma_f$，则 $D=\sigma_f$，进一步由
$D=\varepsilon E$ 得 $E=\sigma_f/\varepsilon$。相对于真空，场强减小 $\varepsilon_r$ 倍，电势差也减小
相同倍数，电容因而增大 $\varepsilon_r$ 倍。这个推导把实验规律、极化图像和宏观场方程连接起来。

不过 $\vect D$ 不能替代 $\vect E$ 描述带电粒子的受力。试探电荷受到的仍是 $q\vect E$，不是
$q\vect D$。在介质界面上，$\vect D$ 的法向跳变由自由面电荷决定，$\vect E$ 的切向分量由静电场保守性
保证连续；只有把两个条件结合起来，才能正确判断电场线如何折转。

\subsection*{线性模型的适用范围}

$\vect P=\varepsilon_0\chi_e\vect E$ 是近似的材料关系。场强过大时，极化可能不再与场成正比；极性分子
的取向还会受到热运动影响，因此介电常数可能随温度变化。交变电场频率很高时，较慢的极化机制来不及响应，
介电常数会随频率变化并出现能量损耗。各向异性晶体中，$\vect P$ 也未必与 $\vect E$ 平行，此时介电常数
应当用张量而不是一个标量表示。

原子极化率 $\alpha$ 是微观量，介电常数是宏观量。若简单地认为每个原子所见局域场就是宏观平均场，便有
$\vect P=N\alpha\vect E$；但邻近偶极子也会改变原子所在位置的场，所以精确关系需要局域场修正。经典的
Clausius--Mossotti 关系写成
\begin{equation}
  \boxed{\frac{\varepsilon_r-1}{\varepsilon_r+2}
  =\frac{N\alpha}{3\varepsilon_0}},
\end{equation}
它展示了如何由宏观介电常数估计微观极化率。该式适用于由近似独立、各向同性分子组成的稀疏介质；对于
强关联固体、极性液体或接近相变的材料，不能机械套用。
